% Gebruiker en Toegangsbeheer
\section{Gebruiker en Toegangsbeheer}
\subsection{Gebruikersgeval}

\begin{figure}[htbp]
    \centering
    \includegraphics[width=\textwidth]{figures/usecase/management.png}
    \caption{Gebruikersgeval vir Gebruiker en Toegangsbeheer}
    \label{fig:gebruiker_toegangsbeheer_use_case}
\end{figure}

Volgens FREQ-WEB-301, FREQ-WEB-302 en FREQ-WEB-303 moet administrasie gebruikers kan toegangsregte en toegangsgeskiedenis van gebruikers sien asook rol en toegangsregte wysig. Hierdie funksies word uitgevoer op die gebruikerswysiging skerm (\autoref{fig:user_edit_screen}).

Verder moet die stelsel funksionaliteit bied om kennisgewings aan gebruikers te stuur (FREQ-WEB-304).

\subsection{Scenario}

\textbf{Akteurs:} Admin\\
\textbf{Voorvereistes:} Admin is ingeteken, Admin het die korrekte regte\\
\textbf{Prekondisie:} Minstens een gebruiker anders as admin bestaan\\
\textbf{Postkondisie:} Gebruiker rol en toegang verander\\
\textbf{Inset:} Gebruiker, gebruiker rolle, blokasie (ja, ne)\\
\textbf{Uitset:} Sukses- of foutboodskap\\
\textbf{Skerms van Toepassing:} 
\begin{itemize}
    \item \autoref{fig:user_edit_screen} (Gebruikerswysiging skerm)
\end{itemize}


\textbf{Hoofvloei:}
\begin{enumerate}
    \item Administrasie gebruiker navigeer na gebruiker bestuur (\autoref{fig:user_edit_screen})
    \item Stelsel toon gebruiker bestuur dashboard (\autoref{fig:user_edit_screen})
    \item Admin soek vir spesifieke gebruiker of blaai deur gebruiker lys
    \item Stelsel toon gebruiker lys met basiese inligting
    \item Admin kies gebruiker om te bestuur (\autoref{fig:user_edit_screen})
    \item Stelsel toon gebruiker besonderhede, rol, en toegangsgeskiedenis (\autoref{fig:user_edit_screen})
    \item Admin kies aksie: rol wysig, toegang blokkeer/deblokkeer, of kyk geskiedenis
    \item Stelsel toon relevante vorm/inligting (\autoref{fig:user_edit_screen})
    \item Admin maak nodige veranderinge
    \item Admin bevestig aksie
    \item Stelsel werk gebruiker inligting by
    \item Stelsel toon bevestiging (\autoref{fig:user_edit_screen})
    \item Gebruikersgeval eindig
\end{enumerate}

\textbf{Alternatiewe vloei A1: Rol wysiging (gewone gebruiker na admin)}
\begin{enumerate}
    \item[7a.] Admin kies "Rol Wysig"
    \item Stelsel toon beskikbare rolle
    \item Admin kies nuwe rol (admin)
    \item Stelsel vra bevestiging vir rol opgradering
    \item Gaan na stap 10
\end{enumerate}

\textbf{Alternatiewe vloei A2: Gebruiker blokkeer/deblokkeer}
\begin{enumerate}
    \item[7b.] Admin kies "Toegang Blokkeer" of "Toegang Deblokkeer"
    \item Stelsel vra rede vir aksie
    \item Admin voer rede in
    \item Gaan na stap 10
\end{enumerate}

\textbf{Alternatiewe vloei A3: Stuur kennisgewings}
\begin{enumerate}
    \item[7c.] Admin kies "Stuur Kennisgewing"
    \item Stelsel bied opsies: individuele gebruiker of stelselwyde
    \item Admin skryf kennisgewingboodskap
    \item Admin kies ontvangers (indien individueel)
    \item Admin stuur kennisgewing
    \item Stelsel stuur kennisgewing aan gekose gebruikers
\end{enumerate}

\textbf{Alternatiewe vloei A4: Kyk toegangsgeskiedenis}
\begin{enumerate}
    \item[7d.] Admin kies "Toegangsgeskiedenis"
    \item Stelsel toon gebruiker se inteken geskiedenis en aktiwiteite
    \item Admin kan periode filter
    \item Gebruikersgeval eindig
\end{enumerate}

\textbf{Uitsonderingsvloei E1: Kan nie eie rol wysig}
\begin{enumerate}
    \item[A1.3.] Admin probeer om sy eie rol te wysig
    \item Stelsel verhoed aksie
    \item Stelsel toon "Kan nie eie rol wysig nie" foutboodskap
    \item Gebruikersgeval eindig
\end{enumerate}
